\documentclass[10pt]{article}
\usepackage[margin=1in]{geometry}
\usepackage{booktabs}
\usepackage{graphicx}
\usepackage{amsmath}
\usepackage{amssymb}
\usepackage[hidelinks]{hyperref}
\usepackage{caption}
\captionsetup{font=small,labelfont=bf}
\usepackage{titlesec}
\titleformat*{\section}{\large\bfseries}
\titleformat*{\subsection}{\normalsize\bfseries}

\newcommand{\stitle}[1]{\vspace{0.4em}\noindent\textbf{#1}\hspace{0.6em}}

\title{Supplementary Material\\[0.3em]
\large Relevant, But Not Useful: Prospective Utility in Agentic Search\\
\normalsize Paper ID: 507 (Anonymous submission)}
\date{}

\begin{document}
\maketitle
\vspace{-2.5em}

\noindent This supplement provides implementation details, the preregistration
and provenance protocol, the agent compliance audit, exact label and feature
definitions, and detailed diagnostic and robustness results referenced in the
main paper.
Section~\ref{sec:agent} details the agent and retrieval stack, including the
stronger relevance selectors;
Section~\ref{sec:candidates} the candidate pool, selector configurations, and
the choice-space consistency audit;
Section~\ref{sec:prereg} the preregistration chain and analysis provenance;
Section~\ref{sec:compliance} the cross-family compliance audit and post-hoc
replication;
Section~\ref{sec:labels} outcome labels and evaluation targets;
Section~\ref{sec:features} observable state, candidate, and continuation
features;
Section~\ref{sec:diagnostic} additional diagnostic results;
Section~\ref{sec:robustness} robustness and sensitivity results; and
Section~\ref{sec:repro} reproducibility information.

\section{Agent and Retrieval Stack}
\label{sec:agent}

\stitle{Search agent.}
The agent follows a ReAct-style loop: at each turn it emits either a search
action (a query wrapped in \texttt{<search>} tags) or a final answer
(wrapped in \texttt{<answer>} tags).
Each search returns the top-3 documents (title and abstract text) from the
fixed corpus, appended to the context as an observation block.
Episodes allow at most 4 search turns; confirmatory continuations use greedy
decoding (temperature $0$) with stop sequences on the action delimiters and on
generation of an observation-opening tag, so that the model cannot hallucinate
its own retrieval results.
Unclosed action tags at end-of-generation are parsed conservatively
(an unclosed \texttt{<search>}/\texttt{<answer>} span is accepted only when it
is unambiguous; occurrences are counted and audited).
A preregistered parser audit found no evidence that these safeguards caused
the observed trajectory differences: parser versions produced 94\% identical
trajectories, compared with 76\% under a same-code control rerun in the same
production environment, and no safeguard-specific mechanism activated in the
confirmatory runs. Additional post-hoc robustness analyses are described in
Secs.~\ref{sec:diagnostic}--\ref{sec:robustness}.

\stitle{Continuation horizon.}
Intervention replays do not receive a fresh search budget. The replayed
episode inherits the number of search calls already consumed by the frozen
prefix, so a state reached after $k$ searches has at most $4-k$ search turns
remaining. On the HotpotQA lockbox intervention arms, the frozen states occur
after 1/2/3/4 consumed searches in 2{,}800/1{,}895/656/136 cases,
respectively. The realized continuation uses 0/1/2/3 additional searches in
2{,}989/1{,}739/608/151 cases, for an average of 0.62 additional search steps.
When the global search budget is exhausted, the agent is instructed to
produce a final answer immediately.

\stitle{Corpus.}
All settings use the official HotpotQA \texttt{enwiki-20171001} abstract
corpus (5{,}233{,}235 documents).
2WikiMultiHopQA and MuSiQue are evaluated against the same fixed corpus;
annotated supporting documents are covered for 92.7\% of 2WikiMultiHopQA and
98.0\% of MuSiQue validation questions, and confirmatory 2WikiMultiHopQA
evaluation is restricted to covered questions (main paper, Sec.~4.1).

\stitle{Retrievers.}
Lexical retrieval uses BM25 (\texttt{bm25s}).
Dense retrieval uses \texttt{intfloat/e5-base-v2} with brute-force inner
product over the full corpus (fp32).
The primary relevance baseline is the MS~MARCO cross-encoder
\texttt{cross-encoder/}\allowbreak\texttt{ms-marco-MiniLM-L6-v2},
scoring each candidate's full title and passage text against the current
query $q_t$ in the original confirmatory configuration.
All retriever configurations are fixed across the five evaluation settings.

\stitle{Stronger relevance selectors.}
We additionally evaluate \texttt{BAAI/bge-reranker-v2-m3} using either the
current query $q_t$ or a state-aware representation containing $q_0$, $q_t$,
and consumed evidence. We also evaluate a Qwen2.5-7B listwise relevance
ranker over the observable state. These models are used only as relevance
selectors; they do not observe intervention outcomes. For BGE rescoring, each
document is represented as its title followed by the first 600 characters of
passage text (tokenizer maximum length 512); the exact full-pool rescoring
audit uses this identical configuration. The MiniLM CE baseline scores the
full passage text (title and text, tokenizer maximum length 384) throughout.

\section{Candidate Pool and Relevance Selectors}
\label{sec:candidates}

For every eligible state, the retrieval pool $\mathcal{C}_t$ is the union of
BM25 top-20 lists under $q_t$, $q_0$, and up to six already-seen document
titles (entity-style queries), plus dense top-20 lists under $q_t$ and $q_0$,
excluding all previously consumed documents
($109.1\pm27.7$ documents per state in the primary evaluation).
Up to eight typed candidates are drawn per state:

\begin{itemize}\itemsep0.15em
  \item \textbf{bm25\_top} --- BM25 top-1 under $q_t$;
  \item \textbf{dense\_top} --- dense top-1 under $q_t$;
  \item \textbf{ce\_top} --- the highest-ranked CE candidate not already
        retained by an earlier standard family; the full-pool CE argmax is
        always retained in $\mathcal{D}_t$ and defines the relevance-selection
        baseline, as detailed below;
  \item \textbf{mid\_rel} --- a BM25 rank 9--16 document under $q_t$;
  \item \textbf{novelty} --- in the pool but outside BM25 top-20 under $q_t$
        and outside dense top-5 (low current-query relevance, reached via
        $q_0$/dense-tail expansion);
  \item \textbf{redundancy} --- high $q_t$ relevance (BM25 top-10) with the
        highest lexical (Jaccard) overlap against already-seen evidence
        (retained only if overlap $>0.05$);
  \item \textbf{entity\_bridge} --- reached only via entity-title queries and
        absent from the BM25 ranking under $q_t$;
  \item \textbf{random} --- a random remaining pool member.
\end{itemize}

\noindent
\textbf{Consistency of the CE baseline.}
Candidates are drawn without replacement, and if the pool-level CE argmax has
already been taken as \textbf{bm25\_top} or \textbf{dense\_top}, it remains in
$\mathcal{D}_t$ under that label while \textbf{ce\_top} takes the next-ranked
document.
Consequently, $\arg\max_{d\in\mathcal{D}_t} R_{\mathrm{CE}}(q_t,d)$ and
$\arg\max_{d\in\mathcal{C}_t} R_{\mathrm{CE}}(q_t,d)$ select the same document,
and the main paper's $d_{\mathrm{CE}}$ (defined over $\mathcal{D}_t$) equals
the pool-level CE selection.
We verified this empirically: in all 3{,}881 evaluated states across the five
settings, the pool-level CE argmax is present in $\mathcal{D}_t$
(zero exceptions).
Standard-retrieval-only analyses restrict the oracle to
\{bm25\_top, dense\_top, ce\_top\}.

\subsection{Choice-Space Consistency Audit}
\label{sec:choicespace}

\stitle{Full-pool argmax retention.}
Table~\ref{tab:s1} reports, for each relevance selector, the fraction of
evaluated states whose full-pool argmax over $\mathcal{C}_t$ is already
contained in the retained intervention set $\mathcal{D}_t$. By construction
this holds for the CE baseline; for the stronger selectors the full-pool
argmax is frequently absent from $\mathcal{D}_t$.

\begin{table}[h]
\centering
\small
\caption{Fraction of states whose full-pool selector argmax is retained in
$\mathcal{D}_t$.}
\label{tab:s1}
\begin{tabular}{lccc}
\toprule
\textbf{Setting} & \textbf{CE} & \textbf{BGE ($q_t$)} & \textbf{BGE-State} \\
\midrule
HotpotQA dev     & 1.000     & 0.668 & 0.366 \\
HotpotQA lockbox & 1.000     & 0.650 & 0.318 \\
2Wiki            & 1.000     & 0.578 & 0.249 \\
Qwen2.5-14B      & 1.000     & 0.706 & 0.372 \\
MuSiQue          & 1.000$^*$ & 0.602 & 0.264 \\
\bottomrule
\end{tabular}
\end{table}
\noindent
$^*$The original logged candidate pools retain the CE full-pool argmax in all
3{,}881 evaluated states. A single MuSiQue mismatch appears only when
reconstructing the pool for the post-hoc audit, due to a dense-retrieval
boundary-rank reconstruction difference; it is not a mismatch in the original
intervention pool.

\stitle{$\mathcal{D}_t$-space relevance quality.}
Table~\ref{tab:s2} reports the support-document hit rate of each selector's
choice within $\mathcal{D}_t$, computed only among states for which
$\mathcal{D}_t$ contains at least one annotated support document.

\begin{table}[h]
\centering
\small
\caption{Support hit rate of selector choices within $\mathcal{D}_t$, computed
only among states for which $\mathcal{D}_t$ contains at least one annotated
support document. The final column reports the BGE-State $-$ CE difference
with question-level 95\% confidence intervals.}
\label{tab:s2}
\begin{tabular}{lcccl}
\toprule
\textbf{Setting} & \textbf{CE} & \textbf{BGE} & \textbf{BGE-State} &
\textbf{BGE-State $-$ CE} \\
\midrule
HotpotQA dev     & 0.753 & 0.846 & 0.900 & +0.147 [0.074, 0.223] \\
HotpotQA lockbox & 0.719 & 0.801 & 0.903 & +0.184 [0.100, 0.268] \\
2Wiki            & 0.707 & 0.796 & 0.790 & +0.082 [0.007, 0.158] \\
Qwen2.5-14B      & 0.749 & 0.819 & 0.852 & +0.103 [0.038, 0.172] \\
MuSiQue          & 0.555 & 0.671 & 0.685 & +0.130 [0.044, 0.216] \\
\bottomrule
\end{tabular}
\end{table}

\stitle{Rollout of missing full-pool argmax candidates.}
Because BGE and BGE-State full-pool argmax candidates are frequently absent
from $\mathcal{D}_t$, we conducted an explicitly post-hoc sensitivity analysis
after the lockbox had been opened. We rolled out every missing full-pool BGE
and BGE-State argmax under the same frozen greedy continuation policy,
yielding 3{,}861 additional intervention arms and no missing outcomes.
Table~\ref{tab:s3} reports the resulting downstream comparison.

\begin{table}[h]
\centering
\small
\caption{Post-hoc full-pool choice-space sensitivity. Entries report
downstream F1 differences between each selector's actual full-pool argmax and
CE selection, with question-level 95\% confidence intervals.}
\label{tab:s3}
\begin{tabular}{lcc}
\toprule
\textbf{Setting} & \textbf{BGE(full)$-$CE} & \textbf{BGE-State(full)$-$CE} \\
\midrule
HotpotQA dev     & $-$.019 [$-$.042, .004] & +.012 [$-$.016, .043] \\
HotpotQA lockbox & +.008 [$-$.010, .026]   & +.024 [$-$.003, .051] \\
2Wiki            & +.010 [$-$.006, .025]   & +.018 [$-$.007, .043] \\
Qwen2.5-14B      & $-$.010 [$-$.029, .010] & $-$.009 [$-$.034, .016] \\
MuSiQue          & $-$.001 [$-$.014, .012] & $-$.000 [$-$.017, .017] \\
\bottomrule
\end{tabular}
\end{table}

\noindent
All 10 confidence intervals include zero. Thus, even when BGE and BGE-State
are allowed to select their actual full-pool argmax candidates, their large
relevance advantage does not yield a statistically detectable downstream gain
over CE. This analysis is post hoc and is reported only as a choice-space
sensitivity check.

\section{Preregistration, Lockbox, and Analysis Provenance}
\label{sec:prereg}

\stitle{Preregistration chain.}
All confirmatory decisions and pre-specified diagnostic analyses were fixed
in a version-controlled design-revisions document \emph{before} the
corresponding data were inspected, including: the primary endpoints and
lockbox analysis freeze list; the replication protocols for
2WikiMultiHopQA, the cross-scale agent, and MuSiQue; and additional
robustness and diagnostic analyses (cross-fitted oracle, question-clustered
inference, pool-robustness, and effect-size decomposition).
No acceptance criterion was relaxed after unblinding at any point; criteria
that failed remain recorded as failures together with the adjudication that
followed.

\stitle{Single-open confirmatory lockbox.}
The held-out HotpotQA lockbox was sealed before predictor development:
question IDs were selected and committed, all models used for confirmatory
evaluation were frozen to disk, and SHA-256 hashes of the frozen artifacts and
design document were recorded before opening. The preregistered confirmatory
lockbox evaluation was executed once. After all confirmatory analyses were
completed, we conducted explicitly post-hoc robustness and diagnostic
analyses on the already opened lockbox. These analyses introduced no new
confirmatory hypothesis testing and are reported separately from the
preregistered results.

\stitle{Confirmatory versus post-hoc analyses.}
Confirmatory (frozen before the corresponding data were inspected): the
primary Oracle--CE endpoint and opportunity rates, the lockbox evaluation,
the representative non-enriched strata, and the split-replay and noise
checks. Explicitly post-hoc and exploratory: the choice-space full-pool
rollout expansion, winner characterization, the within-state
outcome-permutation analysis, budgeted lookahead diagnostics, the
post-rollout continuation scorer, sampled-decoding sensitivity, the Phi-4
cross-family replication, the forced-answer myopic-utility analysis, and
alternative control-matching specifications.

\stitle{Measurement-integrity control.}
A preregistered harness audit compared agent behavior across parser
versions. The strict no-op criterion (99\% identical trajectories) was
\emph{not} met (94\%); a preregistered same-code control (re-running the
identical legacy code twice under production conditions) yielded 76\%
identical trajectories, attributing the discrepancy to serving-level
nondeterminism, and zero activations of the new parsing mechanisms were
observed in confirmatory runs. We therefore report ``no evidence of a
behavioral effect'' rather than a proven no-op, and record the original
criterion as failed.

\section{Cross-Family Compliance Audit}
\label{sec:compliance}

Cross-family agent replications were subject to preregistered
tool-protocol compliance criteria evaluated on a 30-question smoke set
before any confirmatory run (malformed action rate $\le5\%$, valid answer
rate, token-limit truncation rate, and zero answer leakage).
Table~\ref{tab:compliance} summarizes the audit.
Only Qwen2.5-14B passed this preregistered audit; Mistral-7B-v0.3 and
Llama-3.1-8B were excluded \emph{before} confirmatory evaluation, and the
preregistration prohibits re-measuring an excluded model. We therefore
retain Qwen2.5-14B as the only preregistered model-transfer evaluation.
After completion of the confirmatory analyses, we additionally conducted the
explicitly post-hoc Phi-4 cross-family replication described below.

\begin{table}[h]
\centering
\small
\caption{Compliance audit for candidate cross-family agents
(30-question preregistered smoke evaluation; criteria fixed in advance).}
\label{tab:compliance}
\begin{tabular}{lcccl}
\toprule
\textbf{Agent} & \textbf{Malformed} & \textbf{Valid answer} &
\textbf{Token-limit} & \textbf{Decision} \\
\midrule
Qwen2.5-7B-Instruct  & (reference) & --- & --- & primary agent \\
Mistral-7B-Instruct-v0.3 & 42.2\% & 100\% & 0\% & excluded \\
Llama-3.1-8B-Instruct & 5.9\% & 96.7\% & 3.3\% & excluded (near boundary) \\
Qwen2.5-14B-Instruct & 0.0\% & 93.3\% & 6.7\% & \textbf{passed} \\
\bottomrule
\end{tabular}
\end{table}

\subsection{Post-hoc Cross-Family Replication with Phi-4}
\label{sec:phi4}

To test whether the observed prospective headroom is specific to the
Qwen2.5 agent family, we conducted an explicitly post-hoc replication with
\texttt{microsoft/phi-4} after completing the confirmatory analyses. Before
examining replication outcomes, we fixed a separate post-hoc
protocol-compliance screen using the same tool-use criteria. On a
50-question development smoke set, Phi-4 achieved a 92.0\% valid-final
rate, 0.0\% malformed-search rate, 2.0\% token-limit truncation rate, and
zero retrieval leakage, satisfying the fixed compliance thresholds. We then
evaluated Phi-4 on a representative, non-enriched HotpotQA sample of 350
questions and 860 frozen states under the same frozen-state intervention and
greedy continuation protocol. The replication yields significant
Oracle--CE headroom of \textbf{+0.081 F1} (95\% CI [0.063, 0.102]). CE also
significantly exceeds uniform random selection by \textbf{+0.048 F1}
[0.029, 0.067]. Stronger relevance does not eliminate the gap: residual
Oracle--BGE($q_t$) headroom is \textbf{+0.071} [0.053, 0.089], and
Oracle--BGE-State headroom is \textbf{+0.078} [0.060, 0.099]. Query-only
BGE improves slightly over CE (+0.011 [0.002, 0.021]), whereas BGE-State
does not (+0.003 [$-$0.010, 0.016]). Thus, stronger relevance can improve
selection while substantial prospective headroom remains. Because this
experiment was conducted after the confirmatory analyses, we treat it
strictly as a post-hoc cross-family replication.

\section{Outcome Labels and Evaluation Targets}
\label{sec:labels}

\stitle{Prospective utility and headroom.}
For an unseen candidate $d$, raw prospective interventional utility is the
change in downstream answer F1 obtained by injecting $d$ into the frozen
state $s_t$ and resuming the same policy. For a relevance selector $R$,
headroom is
\begin{equation}
H_R(s_t) = Y(d_{\mathrm{PIU}}) - Y(d_R),
\label{eq:headroom}
\end{equation}
the difference between the best measured candidate and the selector's choice
within the same intervention set $\mathcal{D}_t$. Two opportunity
definitions are used and always reported separately: Opp.-Base, the fraction
of states with at least one candidate whose raw intervention outcome
strictly exceeds the base (no-intervention) continuation outcome; and
Opp.-CE, the fraction of states with at least one candidate whose outcome
strictly exceeds that of the CE-selected candidate. The substantive
threshold $H\ge0.25$ F1 distinguishes answer-level changes from marginal
fluctuations.

\stitle{Prospective state-identification target.}
The RQ3 gate predicts, before any intervention outcome is observed, the
binary state-level label
\begin{equation}
Z(s_t) = \mathbb{1}[H_{\mathrm{CE}}(s_t) > 0],
\label{eq:zlabel}
\end{equation}
i.e., whether the current state contains any candidate that would improve on
the CE selection. Gates are evaluated by AUC and average precision against
this retrospectively computed label; the label itself is never available to
the gate.

\stitle{Post-rollout scorer labels.}
The exploratory continuation scorer is trained on candidate pairs whose
preference labels are constructed from the raw downstream F1 outcomes of
their continuations. Pairs whose outcome difference falls inside the
preregistered tie zone ($\tau=0.1$) are excluded from training; this is the
same tie zone used for pairwise outcome labels. F1 is used only to construct
training labels and never appears among inference-time features.

\stitle{Myopic utility.}
For the post-hoc forced-answer analysis, myopic utility is defined
analogously to PIU except that the continuation policy disables further
retrieval and requires an immediate answer after candidate injection. The
resulting outcome measures the candidate's immediate answer-support value
under the frozen state while preventing any contribution from subsequent
retrieval. This makes the utility label explicitly dependent on the
continuation policy: the same candidate can carry different utility under
the full policy and under the forced-answer policy.

\section{Observable State, Candidate, and Continuation Features}
\label{sec:features}

\stitle{Pre-intervention features.}
The interpretable feature vector $\varphi(s_t,d)$ asks what the candidate
\emph{adds} relative to current knowledge.
Content tokens are lowercased alphanumeric tokens with a fixed stopword list
removed; entity proxies are capitalized tokens and multi-digit numbers.
With $D$ the candidate's content-token set, $E_1,\dots,E_k$ the per-document
evidence token sets ($E=\bigcup_i E_i$), $Q$ and $Q_t$ the token sets of the
original question and current query:

\begin{itemize}\itemsep0.15em
  \item \textbf{novelty} $=|D\setminus E|/|D|$ --- fraction of candidate
        content unseen in evidence;
  \item \textbf{redundancy} $=\max_i \mathrm{Jaccard}(D,E_i)$ --- maximal
        overlap with any single consumed document;
  \item \textbf{entity\_gain} --- new entity-proxy types per 100 tokens
        (scale-clipped);
  \item \textbf{query\_align} $=|D\cap Q_t|/|Q_t|$;
  \item \textbf{q0\_align} $=|D\cap Q|/|Q|$;
  \item \textbf{gap\_coverage} $=|D\cap(Q\setminus E)|/|Q\setminus E|$ ---
        coverage of question tokens not yet present in evidence;
  \item \textbf{doc\_len} $=\log(1+\text{token count})$, normalized;
  \item \textbf{turn\_idx}, \textbf{n\_evidence} --- search step and number of
        consumed documents, normalized.
\end{itemize}

\noindent
Semantic variants of the alignment and novelty features (E5 embedding based)
are additionally used in the gate analyses (Sec.~\ref{sec:diagnostic}).
These features are used descriptively in the winner characterization and as
inputs to the prospective gates; all of them are computable before any
intervention outcome is observed.

\stitle{Post-rollout continuation features.}
The exploratory continuation scorer additionally observes features of the
realized continuation of each candidate: normalized CE score of the
candidate; the number of additional searches in the continuation;
termination indicators; trajectory length in tokens and turns; answer length
and an empty-answer indicator; lexical support of the answer from the
injected document; lexical support from subsequently retrieved evidence;
agreement of the continuation answer with the base answer; and an indicator
of whether the answer matches the modal answer among the top-3 CE
candidates' continuations. These features require running the candidate's
continuation and are therefore available only to the post-rollout scorer,
never to the pre-intervention gates.

\section{Additional Diagnostic Results}
\label{sec:diagnostic}

\subsection{Where Utility Winners Appear in the Relevance Ranking}
\label{sec:winnerrank}

Table~\ref{tab:winnerrank} characterizes, among CE-relative opportunity
states, the CE-rank position of utility-maximizing candidates. Winners
concentrate near, but not at, the top of the relevance ranking.

\stitle{Tie handling for winner-rank characterization.}
When multiple candidates attain the same maximal downstream outcome within a
state, we fractionally distribute that state's winner mass across all tied
maximizers, analogously to the candidate-family attribution analysis. Top-$k$
shares are therefore computed from tie-weighted winner mass, and median ranks
from the corresponding weighted empirical distribution.

\begin{table}[h]
\centering
\small
\caption{Tie-aware CE-rank position of utility-maximizing candidates among
CE-relative opportunity states. When multiple candidates attain the same
maximal downstream outcome, each receives equal fractional winner mass within
the state. Top-2/Top-3 report tie-weighted winner mass within CE rank 2/3;
median ranks are computed from the corresponding weighted empirical
distribution, and the final column reports winner mass at CE rank $\ge5$.}
\label{tab:winnerrank}
\begin{tabular}{lccccc}
\toprule
\textbf{Setting} & \textbf{Opp. states} & \textbf{Median CE rank} &
\textbf{Top-2} & \textbf{Top-3} & \textbf{Rank 5+} \\
\midrule
HotpotQA dev     & 81  & 4 & 17.3\% & 38.3\% & 34 \\
HotpotQA lockbox & 129 & 3 & 33.3\% & 50.4\% & 42 \\
2Wiki            & 164 & 4 & 23.8\% & 47.0\% & 58 \\
Qwen2.5-14B      & 140 & 4 & 30.7\% & 45.7\% & 54 \\
MuSiQue          & 187 & 4 & 17.1\% & 35.3\% & 84 \\
\bottomrule
\end{tabular}
\end{table}

\stitle{Outcome-permutation characterization.}
To describe which observable properties distinguish winning candidates, we
compare each feature's observed winner statistic against a within-state
outcome-permutation null. This analysis is exploratory and descriptive; the
reported directions are associations under the tested null, not causal
effects. Table~\ref{tab:permnull} summarizes how many of the five settings
fall outside the permutation null for each feature.

\begin{table}[h]
\centering
\small
\caption{Within-state outcome-permutation characterization of winning
candidates (exploratory). ``Direction'' describes the observed tendency in
the settings outside the null.}
\label{tab:permnull}
\begin{tabular}{lcl}
\toprule
\textbf{Feature} & \textbf{Settings outside null} & \textbf{Direction} \\
\midrule
Document length            & 5/5 & higher \\
Current-query alignment    & 4/5 & higher \\
Unresolved-term coverage   & 4/5 & higher \\
$q_0$ alignment            & 3/5 & higher \\
Redundancy                 & 3/5 & lower \\
Novelty                    & 2/5 & mixed / limited \\
Entity gain                & 0/5 & no consistent signal \\
\bottomrule
\end{tabular}
\end{table}

\subsection{Detailed Prospective State-Identification Results}
\label{sec:gatedetail}

The RQ3 gate predicts $Z(s_t)=\mathbb{1}[H_{\mathrm{CE}}(s_t)>0]$ from
pre-intervention information only: CE margin and CE score statistics,
candidate-aggregate features, and probe signals. Table~\ref{tab:gateauc}
reports discrimination of the full shallow gate; across post-hoc gate
families, the best per-setting AUC ranges from 0.567 to 0.637, and CE margin
or CE confidence alone performs approximately at chance.

\begin{table}[h]
\centering
\small
\caption{Prospective state-identification AUC of the full shallow gate
against $Z(s_t)$, by evaluation setting. Across post-hoc gate families, the
best per-setting AUC ranges from 0.567 to 0.637; for Qwen2.5-14B the best
family is the candidate-aggregate gate (0.568) rather than the full gate.}
\label{tab:gateauc}
\begin{tabular}{lc}
\toprule
\textbf{Setting} & \textbf{Full shallow gate AUC} \\
\midrule
HotpotQA dev     & 0.599 \\
HotpotQA lockbox & 0.637 \\
2Wiki            & 0.527 \\
Qwen2.5-14B      & 0.560 \\
MuSiQue          & 0.505 \\
\bottomrule
\end{tabular}
\end{table}

\begin{table}[h]
\centering
\small
\caption{Gate-family comparison on the HotpotQA lockbox (prevalence of
$Z(s_t)=1$: 0.188). The full shallow gate attains an average-precision lift
of 1.49 relative to prevalence.}
\label{tab:gatefamily}
\begin{tabular}{lc}
\toprule
\textbf{Gate family} & \textbf{Lockbox AUC} \\
\midrule
CE margin                    & 0.479 \\
CE score statistics          & 0.494 \\
Candidate aggregates ($\varphi$) & 0.547 \\
MiniLM (fine-tuned)          & 0.500 \\
DeBERTa (fine-tuned)         & 0.525 \\
Full shallow gate            & 0.637 \\
\bottomrule
\end{tabular}
\end{table}

\subsection{Budgeted Lookahead and Simple Aggregation}
\label{sec:lookahead}

Table~\ref{tab:lookahead} reports the fraction of oracle headroom recovered
when the intervention budget is limited to the top-$k$ CE-ranked candidates.

\begin{table}[h]
\centering
\small
\caption{Fraction of Oracle--CE headroom recovered under a top-$k$
CE-ranked intervention budget.}
\label{tab:lookahead}
\begin{tabular}{lcc}
\toprule
\textbf{Setting} & \textbf{Oracle@2 recovery} & \textbf{Oracle@3 recovery} \\
\midrule
HotpotQA dev     & 30.9\% & 64.0\% \\
HotpotQA lockbox & 46.4\% & 65.4\% \\
2Wiki            & 37.3\% & 60.8\% \\
Qwen2.5-14B      & 46.3\% & 70.6\% \\
MuSiQue          & 32.2\% & 54.2\% \\
\bottomrule
\end{tabular}
\end{table}

\noindent
Recovery continues to rise with budget: 68.2--80.0\% of oracle headroom at
$k=4$ and 85.8--95.2\% at $k=6$ across the five settings. Simple aggregation
over the rolled-out candidates does not convert this budgeted signal into a
selector: majority voting over top-$k$ continuation answers degrades
performance relative to CE selection in every setting (negative recovery,
$-5.9\%$ to $-58.0\%$ at $k\in\{4,6\}$), and early-termination heuristics are
inconsistent across settings. One reason is answer modality: the useful
oracle candidate's continuation answer matches the modal answer among top-3
continuations in only 22.5--28.7\% of cases, whereas the CE candidate's
answer is modal in 34.1--47.9\%; useful answers are disproportionately
non-modal.

\subsection{Exploratory Post-Rollout Continuation Scorer}
\label{sec:scorer}

An unconstrained selector was first evaluated and did not improve over CE on
HotpotQA development (development gain: $-0.022$ F1). We therefore treat
this variant as unsuccessful rather than selecting it for transfer
evaluation. The reported variant is CE-sticky: it overrides the CE choice
only when the scorer margin exceeds a threshold selected on HotpotQA
development and then frozen for all non-development evaluations. This is a
post-rollout reranker requiring up to $k$ continuation rollouts per state;
it is not a pre-intervention selector, and its results are explicitly
exploratory.

\begin{table}[h]
\centering
\small
\caption{CE-sticky post-rollout continuation scorer (exploratory).
Overrides counts states where the scorer replaces the CE choice.}
\label{tab:scorer}
\begin{tabular}{lcccc}
\toprule
\textbf{Setting} & \textbf{Overrides} & \textbf{Gain vs. CE} &
\textbf{95\% CI} & \textbf{Oracle@3 recovered} \\
\midrule
HotpotQA dev     & 12/579 & +0.0011 & [$-$0.003, 0.005]  & 2.1\% \\
HotpotQA lockbox & 13/686 & +0.0035 & [0.0007, 0.0071]   & 4.6\% \\
2Wiki            & 39/819 & +0.0075 & [0.0017, 0.0143]   & 8.7\% \\
Qwen2.5-14B      & 49/838 & +0.0096 & [0.0025, 0.0172]   & 14.0\% \\
MuSiQue          & 58/959 & +0.0004 & [$-$0.004, 0.005]  & 0.8\% \\
\bottomrule
\end{tabular}
\end{table}

\noindent
The continuation scorer is fit only on the designated training split; the
CE-sticky margin is selected on HotpotQA development and then frozen for all
non-development evaluations.

\subsection{Trajectory Decomposition}
\label{sec:trajdecomp}

The main paper reports the lockbox decomposition in detail together with the
query-change stratum for the remaining settings; the full per-setting table
(strata: no search / all / any search / both search / query changed, with
state counts, headroom confidence intervals, and Opp.-CE) is provided with
the frozen analysis reports in the repository. On the lockbox, the
query-changed stratum ($n=175$) shows headroom of +0.311 [0.244, 0.380] with
Opp.-CE of 44.6\%. The decomposition is computed over the post-intervention
behavior of the two compared arms ($d_{\mathrm{PIU}}$ and $d_{\mathrm{CE}}$);
its denominator therefore differs from the all-arm average of 0.62
additional search steps reported in Sec.~\ref{sec:agent}. Post-intervention
strata are descriptive and should not be interpreted causally.

\subsection{Myopic Utility versus Full Prospective Utility}
\label{sec:myopic}

The trajectory decomposition above is descriptive because its strata are
defined after intervention. We therefore conduct a complementary
forced-answer analysis on the HotpotQA lockbox. For each frozen state and
candidate, the full-PIU arm injects the candidate and resumes the normal
greedy continuation policy, whereas the myopic arm injects the same
candidate but disables further retrieval and requires an immediate answer.
All other state, candidate, and evaluation conditions are held fixed.
Retrieval-free myopic utility is moderately correlated with full PIU
(candidate-level Spearman $\rho=0.652$), but its within-state ordering is
much coarser: mean within-state Spearman correlation is 0.392 and the
myopic maximum is tied in 84.1\% of states. Among CE-relative opportunity
states, strict winner agreement with full PIU is only 27.1\%. The myopic
oracle still exhibits substantial headroom (+0.083 F1, 95\% CI
[0.063, 0.104]), showing that immediate answer effects are real. However,
when ties are resolved in favor of CE, myopic selection gains only +0.015 F1
[0.003, 0.027] on full-continuation outcomes, recovering 12.7\% of the full
Oracle--CE headroom. Thus, immediate answer improvement alone is
insufficient to identify which evidence is most valuable under continued
search.

\begin{table}[h]
\centering
\small
\caption{Post-hoc forced-answer comparison of myopic utility and full PIU on
the HotpotQA lockbox.}
\label{tab:myopic}
\begin{tabular}{lc}
\toprule
\textbf{Quantity} & \textbf{Result} \\
\midrule
Candidate-level Spearman $\rho$ & .652 \\
Mean within-state Spearman & .392 \\
States with tied myopic maximum & 84.1\% \\
Strict winner agreement (opportunity states) & 27.1\% \\
Full Oracle--CE & +.117 [.092,.142] \\
Myopic Oracle--CE & +.083 [.063,.104] \\
Myopic selection on full outcomes & +.015 [.003,.027] \\
Full headroom recovered & 12.7\% \\
\bottomrule
\end{tabular}
\end{table}

\section{Robustness and Sensitivity Results}
\label{sec:robustness}

\stitle{Exact random selection and representative strata.}
The random baseline is computed in closed form as the exact expectation of a
uniform draw over $\mathcal{D}_t$ (the uniform average of the candidate
outcomes), and compared with CE selection using paired question-level
confidence intervals. CE selection is significantly better than this exact
random baseline in all five settings (CE $-$ Random from +0.0228 to +0.0720
F1), so the reported headroom is measured beyond a genuinely informative
baseline. On the representative non-enriched strata, the Oracle--CE effect
persists: question-level estimates range from +0.062 to +0.159 F1 across the
five settings, with every confidence interval excluding zero. Failure
enrichment therefore does not create the Oracle--CE effect; it only
concentrates evaluation on states where selection matters.

\stitle{Opportunity magnitude.}
Table~\ref{tab:oppmag} decomposes both opportunity definitions by effect
threshold. By definition, Opp.-Base $\ge$ Opp.-CE within every row, and
$>0$ rates bound the $\ge.25$ rates within each definition.

\begin{table}[h]
\centering
\small
\caption{Opportunity rates under both definitions and effect thresholds.}
\label{tab:oppmag}
\begin{tabular}{lcccc}
\toprule
\textbf{Setting} & \textbf{Opp.-Base $>0$} & \textbf{Opp.-Base $\ge.25$} &
\textbf{Opp.-CE $>0$} & \textbf{Opp.-CE $\ge.25$} \\
\midrule
HotpotQA dev     & 30.1\% & 26.4\% & 14.0\% & 12.4\% \\
HotpotQA lockbox & 29.3\% & 27.8\% & 18.8\% & 16.9\% \\
2Wiki            & 26.4\% & 24.8\% & 20.0\% & 17.6\% \\
Qwen2.5-14B      & 32.1\% & 29.1\% & 16.7\% & 14.3\% \\
MuSiQue          & 21.8\% & 18.8\% & 19.5\% & 16.1\% \\
\bottomrule
\end{tabular}
\end{table}

\stitle{Noise, replay, and cross-fitted robustness.}
Label reliability on the replay-audited subset is 97.8--99.0\% (three
greedy continuations per intervention on 20\% of states; 115/579 primary
states and 137/686 lockbox states). A fresh 150-question cross-fitted
experiment separates candidate selection from outcome evaluation using two
independent rollout sets, yielding an Oracle--CE estimate of +0.111 F1
(95\% CI [0.080, 0.144]), nearly identical to the naive +0.114 estimate on
the same sample (A$\to$B: +0.117; B$\to$A: +0.104).

\stitle{Post-hoc sampled-decoding policy sensitivity.}
To avoid selecting and evaluating the oracle candidate on the same
stochastic realization, we use held-out replicate evaluation. For each
fold, candidate selection is performed using the remaining sampled
replicates and the selected candidate is evaluated on the held-out
replicate. Across the full 686-state lockbox, the three held-out
Oracle--CE estimates are +0.0298, +0.0281, and +0.0407 F1. Thus, realized
PIU remains positive under sampled continuation but is substantially
smaller than under the primary greedy policy, demonstrating policy
dependence. The sampled-decoding runs use $T=0.7$, top-$p=1.0$, with fixed
seeds 101, 102, and 103; prompts, stop sequences, parsing, retrieval, and
search-budget logic are otherwise unchanged.

\stitle{Matched intervention controls.}
The selection-specific controlled contrast is the primary control estimand:
for the candidates actually selected by CE and by the PIU oracle,
Table~\ref{tab:matchedcontrol} reports each selection's outcome relative to
its matched control and the resulting controlled Oracle--CE contrast.

\begin{table}[h]
\centering
\small
\caption{Selection-specific matched-control contrasts (F1).}
\label{tab:matchedcontrol}
\begin{tabular}{lccc}
\toprule
\textbf{Setting} & \textbf{CE $-$ control} & \textbf{Oracle $-$ control} &
\textbf{Controlled Oracle$-$CE} \\
\midrule
HotpotQA dev     & +0.085 & +0.160 & +0.075 \\
HotpotQA lockbox & +0.044 & +0.152 & +0.109 \\
2Wiki            & +0.025 & +0.164 & +0.138 \\
Qwen2.5-14B      & +0.066 & +0.162 & +0.097 \\
MuSiQue          & +0.030 & +0.113 & +0.083 \\
\bottomrule
\end{tabular}
\end{table}

\noindent
The paired controlled Oracle$-$CE contrast excludes zero in all five
settings (e.g., HotpotQA dev +0.0748 [0.0516, 0.0992]; 2Wiki +0.1383
[0.1103, 0.1678]). The near-zero all-arm average is not the primary
selection-specific estimand: many intervention candidates are never
selected by either CE or the oracle. For the candidates actually selected by
CE and by the PIU oracle, the controlled contrast remains positive in all
five settings. On the lockbox, the all-arm controlled average is +0.0060 F1
with a 95\% CI of [$-$0.0017, 0.0137]; excluding controls already present in
the agent context changes it only slightly (+0.0060 to +0.0051; exclusion CI
[$-$0.0027, 0.0127]), and for interventions assigned a second independently
sampled control, effect signs agree in 87--91\% of cases.

\stitle{Additional matching-specification sensitivity.}
As a post-hoc exploratory robustness analysis, we reran all 5{,}487 lockbox
interventions using three relaxed control-matching specifications: removing
the BM25 rank-bucket constraint, removing the document-length constraint, or
removing both. The resulting all-arm controlled estimates are +0.0099,
+0.0089, and +0.0086 F1, respectively, compared with +0.0060 under the full
confirmatory specification (Table~\ref{tab:matchspec}); relaxing the
matching constraints therefore does not attenuate the diagnostic.
Table~\ref{tab:matchcheck} verifies that the relaxed specifications actually
change the sampled controls.

\begin{table}[h]
\centering
\small
\caption{Post-hoc sensitivity to matched-control specifications on the
HotpotQA lockbox. Estimates are question-level means with question-bootstrap
95\% confidence intervals.}
\label{tab:matchspec}
\begin{tabular}{lcc}
\toprule
\textbf{Matching specification} & \textbf{Mean $\Delta$F1} & \textbf{95\% CI} \\
\midrule
Full confirmatory & +0.0060 & [$-$0.0017, 0.0137] \\
No rank bucket    & +0.0099 & [0.0024, 0.0177] \\
No length         & +0.0089 & [0.0009, 0.0171] \\
Neither           & +0.0086 & [0.0006, 0.0167] \\
\bottomrule
\end{tabular}
\end{table}

\begin{table}[h]
\centering
\small
\caption{Manipulation check for relaxed control-matching specifications.}
\label{tab:matchcheck}
\begin{tabular}{lccc}
\toprule
\textbf{Specification} & \textbf{Rank-bucket match} &
\textbf{Length within $\pm25\%$} & \textbf{Fallback} \\
\midrule
No rank bucket & 11.0\% & 97.0\% & 0.3\% \\
No length      & 59.4\% & 22.3\% & 1.8\% \\
Neither        & 9.8\%  & 20.6\% & 0.1\% \\
\bottomrule
\end{tabular}
\end{table}

\stitle{Oracle tie handling.}
The intervention pipeline resolves equal maxima by the fixed candidate
enumeration order bm25\_top $\to$ dense\_top $\to$ ce\_top $\to$ mid\_rel
$\to$ novelty $\to$ redundancy $\to$ entity\_bridge $\to$ random. On the
HotpotQA lockbox, 580 of 686 states have multiple candidates tied at the
maximum raw utility; most of these ties occur because the maximum utility is
zero. Among states with a positive maximum, 100 states (14.6\% of all
lockbox states) contain a tied maximizer. This tie rule cannot change
Oracle--CE headroom because tied candidates have identical downstream F1
within the state, but it can affect descriptive candidate attribution and
winner-rank statistics. We therefore use fractional tie weighting for both
candidate-family attribution and the winner-rank characterization
(Sec.~\ref{sec:winnerrank}). For this reason, family shares below divide
oracle-winning mass fractionally across tied maximizers.

\stitle{Candidate-family attribution.}
Because multiple candidates can attain the same maximal downstream outcome,
assigning each state to the first enumerated maximizer biases family
counts. We therefore fractionally divide each state's oracle-winning mass
across all tied maximizers. Under this tie-split attribution, the three
standard-retrieval candidate families (bm25\_top, dense\_top, and ce\_top)
account for 56.8\% of oracle-winning mass, while the five auxiliary
families account for 43.2\%. Table~\ref{tab:family} reports the per-family
breakdown, and Table~\ref{tab:poolsens} extends the pool-sensitivity
analysis to all five evaluation settings.

\begin{table}[h]
\centering
\small
\caption{Candidate-family attribution on the HotpotQA lockbox (686 states;
201 states with a positive opportunity). Oracle-winning mass is fractionally
divided across candidates tied at the maximum downstream outcome. The final
column applies the pipeline's fixed first-maximum tie rule and is shown only
to illustrate enumeration bias.}
\label{tab:family}
\begin{tabular}{lccc}
\toprule
\textbf{Family} & \textbf{Mean $U$} & \textbf{Tie-split share} &
\textbf{First-max share} \\
\midrule
bm25\_top      & +0.0180 & 8.9\%  & 25.4\% \\
dense\_top     & +0.0638 & 25.7\% & 33.3\% \\
ce\_top        & +0.0652 & 22.2\% & 21.9\% \\
mid\_rel       & +0.0182 & 7.1\%  & 3.0\% \\
novelty        & +0.0239 & 8.7\%  & 6.0\% \\
redundancy     & +0.0172 & 9.0\%  & 4.0\% \\
entity\_bridge & +0.0210 & 9.4\%  & 4.5\% \\
random         & +0.0255 & 8.9\%  & 2.0\% \\
Standard retrieval & --- & 56.8\% & --- \\
Auxiliary families & --- & 43.2\% & --- \\
\bottomrule
\end{tabular}
\end{table}

\begin{table}[h]
\centering
\small
\caption{Candidate-pool sensitivity across all five evaluation settings.
Standard and auxiliary columns show tie-split oracle-winning mass.
Standard-only headroom restricts the oracle to BM25, dense-retrieval, and CE
candidates. Bracketed values are question-level 95\% confidence intervals.}
\label{tab:poolsens}
\begin{tabular}{lcccc}
\toprule
\textbf{Setting} & \textbf{Standard} & \textbf{Aux.} &
\textbf{Full headroom} & \textbf{Standard-only} \\
\midrule
HotpotQA dev     & 60.0\% & 40.0\% & +0.082 & +0.042 [0.029, 0.058] \\
HotpotQA lockbox & 56.8\% & 43.2\% & +0.117 & +0.078 [0.059, 0.100] \\
2Wiki            & 50.4\% & 49.6\% & +0.142 & +0.083 [0.064, 0.104] \\
Qwen2.5-14B      & 61.8\% & 38.2\% & +0.097 & +0.071 [0.053, 0.091] \\
MuSiQue          & 42.0\% & 58.0\% & +0.097 & +0.047 [0.035, 0.060] \\
\bottomrule
\end{tabular}
\end{table}

\noindent
Restricting the intervention pool to standard retrieval candidates leaves
positive Oracle--CE headroom in all five settings, with every
question-level 95\% confidence interval excluding zero. Auxiliary
candidates are not inert: they account for 38--58\% of tie-split
oracle-winning mass across settings, with the largest share on MuSiQue.

\stitle{Cost accounting.}
96.2\% of rollout token cost is prompt tokens (context replay), implying
that amortized approximations of prospective intervention should target
context reuse rather than generation length.

\section{Reproducibility}
\label{sec:repro}

\stitle{Software and hardware.}
vLLM 0.10.2, PyTorch 2.8.0 (CUDA 12.8), Transformers 4.55.2,
\texttt{bm25s} 0.3.10, scikit-learn 1.9.0, Python 3.11; NVIDIA H100 80GB.
Confirmatory evaluations use greedy decoding; fixed seeds are used for
sampling, stratification, and control matching. The post-hoc
sampled-decoding sensitivity analysis uses $T=0.7$, top-$p=1.0$, with fixed
seeds 101, 102, and 103.

\stitle{Scale of the study.}
Per-setting confirmatory sizes (questions\,/\,frozen states): HotpotQA dev
297\,/\,579; lockbox 350\,/\,686; 2Wiki 350\,/\,819; Qwen2.5-14B
350\,/\,838; MuSiQue 349\,/\,959. The cross-fitted oracle uses 150 fresh
questions with two independent rollout sets. The post-hoc choice-space
expansion adds 3{,}861 full-pool BGE/BGE-State intervention arms. The
sampled-decoding sensitivity evaluates 6{,}173 lockbox base/intervention
arms under three independent stochastic replicates. The post-hoc Phi-4
cross-family replication uses a representative, non-enriched 350-question
HotpotQA sample containing 860 frozen states and 6{,}870 intervention arms.
The forced-answer myopic analysis evaluates 6{,}173 lockbox arms (686 base
continuations and 5{,}487 candidate interventions), with retrieval disabled
after candidate injection.

\stitle{Analysis provenance.}
The continuation scorer is fit only on the designated training split; the
CE-sticky margin is selected on HotpotQA development and then frozen for
all non-development evaluations. All post-hoc analyses (choice-space
expansion, winner characterization, lookahead, scorer, sampled decoding,
Phi-4 cross-family replication, myopic-utility analysis, and matching
sensitivity) are labeled as such in the frozen reports.

\stitle{Code and records.}
Anonymized code, the preregistration document, and frozen analysis reports
are available at the repository linked in the abstract. All confirmatory
numbers in the main paper trace to a preregistered criterion and a frozen
report file in that repository.

\end{document}
